%% Essai de traction : courbe contrainte-déformation d'un acier doux.
%% Partie linéaire (domaine élastique, pente E), limite élastique Re,
%% palier plastique, limite à la rupture Rr, puis rupture.
\documentclass[tikz,border=4pt]{standalone}
\usepackage{liaisons}
\begin{document}
\begin{tikzpicture}[x=1cm,y=1cm,
  axe/.style={-{Stealth[length=5pt]}, line width=0.7pt},
  courbe/.style={line width=1.6pt, solideB, line cap=round, line join=round},
  rappel/.style={line width=0.5pt, dash pattern=on 3pt off 2pt, black!60},
  cote/.style={{Stealth[length=4pt]}-{Stealth[length=4pt]}, line width=0.7pt, solideD}]

% --- axes --------------------------------------------------------------
\draw[axe] (0,0) -- (0,3.75);
\draw[axe] (0,0) -- (5.3,0);
\node[font=\small, anchor=north east, inner sep=2pt] at (0,0) {$0$};
\node[font=\small, anchor=north east, inner sep=3pt] at (5.35,-0.1)
     {Déformation $\varepsilon = \Delta L / L$};
\node[font=\small, anchor=south west, align=left, inner sep=2pt] at (-0.35,3.8)
     {Contrainte $\sigma = F / S$};

% --- repères Re et Rr ---------------------------------------------------
\draw[rappel] (0,2.2) -- (1.0,2.2);
\draw[rappel] (0,3.15) -- (3.5,3.15) (3.5,3.15) -- (3.5,0);
\node[font=\small, anchor=east, inner sep=3pt] at (0,2.2) {$R_e$};
\node[font=\small, anchor=east, inner sep=3pt] at (0,3.15) {$R_r$};

% --- la courbe ---------------------------------------------------------
\draw[courbe] (0,0) -- (1.0,2.2)
  .. controls (1.08,2.42) and (1.18,2.0) .. (1.32,1.98)
  .. controls (1.55,1.95) and (1.62,2.06) .. (1.85,2.14)
  .. controls (2.55,2.62) and (2.95,3.05) .. (3.5,3.15)
  .. controls (3.95,3.22) and (4.25,3.0) .. (4.55,2.62);
\draw[line width=1.2pt, solideA] (4.45,2.72) -- (4.72,2.5) (4.45,2.5) -- (4.72,2.72);
\node[font=\small, text=solideA, anchor=west, inner sep=3pt] at (4.7,2.5) {Rupture};

% --- module de Young : pente de la partie linéaire ---------------------
\draw[line width=0.6pt, black!60] (0.34,0.75) -- (0.68,0.75) -- (0.68,1.5);
\node[font=\small, anchor=north west, inner sep=1pt] at (0.7,1.42) {$E$};
\node[font=\small, anchor=south west, inner sep=2pt, rotate=65] at (0.22,0.5) {$\sigma = E \cdot \varepsilon$};

% --- zones élastique et plastique --------------------------------------
\draw[cote] (0.03,-0.95) -- (1.0,-0.95);
\draw[cote] (1.0,-0.95) -- (4.6,-0.95);
\draw[rappel] (1.0,2.2) -- (1.0,-1.05);
\draw[rappel] (4.6,-0.75) -- (4.6,-1.05);
\draw[line width=0.4pt, black!45] (0.5,-1.0) -- (0.5,-1.5);
\node[font=\small, text=solideD, anchor=north, inner sep=3pt] at (0.5,-1.5) {Zone élastique};
\node[font=\small, text=solideD, anchor=north, inner sep=3pt] at (2.85,-1.0) {Zone de déformation plastique};
\end{tikzpicture}
\end{document}
